% Brüche und Dezimalzahlen – bank/brueche-dezimalzahlen/ (zusammenbau v0.4, Heft)
\einheitenkopf[t2]{Brüche und Dezimalzahlen}
\setcounter{aufgabe}{3}

\begin{aufgabe}{Anteil bestimmen – Prüfungsaufgaben}
\begin{teile}
\teil Markiert sind die Sektoren A und E. Welcher Anteil des Kreises ist markiert? \hfill (P10 2019 OS)\\ \kreuz{$\frac{2}{5}$}\\ \kreuz{$\frac{3}{9}$}\\ \kreuz{$\frac{2}{9}$}\\ \kreuz{$\frac{3}{5}$}

\kreisdiagramm{A/80,B/80,C/80,D/80,E/40}
\teil Markiert sind die Sektoren A und B. Welcher Anteil des Kreises ist markiert? \hfill (P10 2019 OS)\\ \kreuz{$\frac{2}{6}$}\\ \kreuz{$\frac{3}{8}$}\\ \kreuz{$\frac{2}{8}$}\\ \kreuz{$\frac{3}{6}$}

\kreisdiagramm{A/90,B/45,C/90,D/45,E/45,F/45}
\teil Das Rechteck besteht aus $12$ Kästchen. Die beiden schrägen Linien bilden mit der unteren Seite ein Dreieck. Welcher Anteil des Rechtecks ist das Dreieck? \hfill (P10 2024 OS) \\

\begin{ksys}[xmin=0,xmax=6,ymin=0,ymax=2,ablesen] \funktionab{\x}{}{0}{2} \funktionab{4-\x}{}{2}{4} \end{ksys}
\leerfeld
\teil Das Quadrat besteht aus $16$ Kästchen. Die beiden schrägen Linien bilden mit der oberen Seite ein Dreieck. Welcher Anteil des Quadrats ist das Dreieck? \hfill (P10 2024 OS) \\

\begin{ksys}[xmin=0,xmax=4,ymin=0,ymax=4,ablesen] \funktionab{4-\x}{}{0}{2} \funktionab{\x}{}{2}{4} \end{ksys}
\leerfeld
\teil Welcher Anteil der Fläche ist grau? Gib ihn als Bruch und in Prozent an. \hfill (P10 2014 OS) \\

\bruchrechteck{13}{25}
Bruch: \leerfeld Prozent: \leerfeld[\%]
\teil Welcher Anteil der Figur ist grau? Gib ihn als Bruch und in Prozent an. \hfill (P10 2014 OS) \\

\bruchrechteck{11}{20}
Bruch: \leerfeld Prozent: \leerfeld[\%]
\teil Markiere $30\,\%$ der Fläche. \hfill (P10 2022 OS)

\bruchrechteck{0}{20}
\teil Markiere $40\,\%$ der Fläche. \hfill (P10 2022 OS)

\bruchrechteck{0}{25}
\teil Markiere $25\,\%$ der Figur. \hfill (P10 2023 OS)

\bruchrechteck{0}{32}
\teil Färbe $25\,\%$ der Figur. \hfill (P10 2023 OS)

\bruchrechteck{0}{36}
\teil Schraffiere $\frac{4}{5}$ der Fläche. \hfill (P10 2017 OS)

\bruchrechteck{0}{20}
\teil Schraffiere $\frac{2}{3}$ der Fläche. \hfill (P10 2017 OS)

\bruchrechteck{0}{18}
\teil Markiere $\frac{5}{12}$ des Rechtecks. \hfill (P10 2021 OS)

\bruchrechteck{0}{12}
\teil Markiere $\frac{4}{9}$ des Rechtecks. \hfill (P10 2021 OS)

\bruchrechteck{0}{9}
\teil Kennzeichne ein Achtel der Fläche des Quadrats. \hfill (P10 2025 OS)

\raute{4}{4}
\teil Kennzeichne drei Viertel der Kreisfläche. \hfill (P10 2025 OS)

\kreisleer[2]
\teil Bei welcher Figur ist genau $\frac{1}{6}$ grau? \hfill (P10 2026 FOR)\\ \kreuz{P}\\ \kreuz{Q}\\ \kreuz{R}\\ \kreuz{S}

P \streifen[0]{20}{}{} \quad Q \kreissektor{60}{} \quad R \bruchrechteck{1}{4} \quad S \bruchkreis{1}{3}
\teil Bei welcher Figur ist genau $\frac{1}{5}$ grau? \hfill (P10 2026 FOR)\\ \kreuz{P}\\ \kreuz{Q}\\ \kreuz{R}\\ \kreuz{S}

P \bruchkreis{1}{4} \quad Q \streifen[0]{30}{}{} \quad R \bruchrechteck{1}{6} \quad S \kreissektor{72}{}
\end{teile}
\end{aufgabe}

\begin{aufgabe}{Bruchteil einer Größe – Prüfungsaufgaben}
\begin{teile}
\teil Ein Aquarium fasst $120$ l und ist zu $\frac{3}{4}$ gefüllt. Wie viele Liter fehlen bis zum Rand? \hfill (P10 2015 OS) \\ \leerfeld[l]
\teil Ein Heizöltank fasst $1\,500$ l und ist zu $\frac{2}{5}$ gefüllt. Wie viele Liter passen noch hinein? \hfill (P10 2015 OS) \\ \leerfeld[l]
\teil $\frac{3}{5}$ von $2{,}5$ kg – wie viel Kilogramm? \hfill (P10 2018 OS) \\ \leerfeld[kg]
\teil Wie viel Gramm sind $\frac{3}{8}$ von $1{,}6$ kg? \hfill (P10 2018 OS) \\ \leerfeld[g]
\end{teile}
\end{aufgabe}

\begin{aufgabe}{Kürzen und Erweitern – Prüfungsaufgaben}
\begin{teile}
\teil Welcher Anteil der Figur ist grau? Gib ihn als vollständig gekürzten Bruch und in Prozent an. \hfill (P10 2014 OS) \\

\bruchrechteck{12}{16}
Bruch: \leerfeld Prozent: \leerfeld[\%]
\teil Welcher Anteil der Fläche ist grau? Kürze vollständig und gib den Anteil auch in Prozent an. \hfill (P10 2014 OS) \\

\bruchrechteck{10}{25}
Bruch: \leerfeld Prozent: \leerfeld[\%]
\end{teile}
\end{aufgabe}

\begin{aufgabe}{Vergleichen – Prüfungsaufgaben}
\begin{teile}
\teil Gib eine Zahl an, die zwischen $\frac{1}{4}$ und $\frac{3}{5}$ liegt. \hfill (P10 2014 OS) \\ \leerfeld
\teil Gib eine Zahl an, die zwischen $\frac{2}{5}$ und $\frac{3}{4}$ liegt. \hfill (P10 2014 OS) \\ \leerfeld
\end{teile}
\end{aufgabe}

\begin{aufgabe}{Umwandeln – Prüfungsaufgaben}
\begin{teile}
\teil Welche Aussage ist wahr? Kreuze an. \hfill (P10 2015 OS)\\ \kreuz{$2{,}5 < \frac{5}{2}$}\\ \kreuz{$\frac{13}{5} > \frac{5}{2}$}\\ \kreuz{$\sqrt{5} > \frac{5}{2}$}
\teil Welche Aussage ist wahr? Kreuze an. \hfill (P10 2015 OS)\\ \kreuz{$1{,}75 < \frac{7}{4}$}\\ \kreuz{$\frac{9}{5} > \frac{7}{4}$}\\ \kreuz{$\sqrt{3} > \frac{7}{4}$}
% AUSGELASSEN brueche-dezimalzahlen-e4-k1-s8-v3: Missing $ inserted.
% \teil Setze $<$, $=$ oder $>$ ein: $2\,\%$ __ $0{,}2$ \hfill (P10 2018 OS) \\ \leerfeld
\teil \textit{(ausgelassen)}
% AUSGELASSEN brueche-dezimalzahlen-e4-k1-s8-v4: LaTeX Error: Command \item invalid in math m
% \teil Setze $<$, $=$ oder $>$ ein: $0{,}3$ __ $3\,\%$ \hfill (P10 2018 OS) \\ \leerfeld
\teil \textit{(ausgelassen)}
\end{teile}
\end{aufgabe}

\begin{aufgabe}{Vergleichen und Ordnen – Prüfungsaufgaben}
\begin{teile}
\teil Welcher Wert ist der kleinste: $5{,}5$; $0{,}55$; $0{,}5^2$; $55\,\%$? \hfill (P10 2023 OS) \\ \leerfeld
\teil Welcher Wert ist der kleinste: $0{,}7^2$; $0{,}5$; $4{,}9$; $45\,\%$? \hfill (P10 2023 OS) \\ \leerfeld
\teil Ordne aufsteigend: $-\frac{3}{4}$; $1{,}7$; $-0{,}749$; $\sqrt{3}$ \hfill (P10 2014 OS) \\ \leerfeld $<$ \leerfeld $<$ \leerfeld $<$ \leerfeld
\teil Ordne aufsteigend: $-\frac{2}{5}$; $2{,}2$; $-0{,}41$; $\sqrt{5}$ \hfill (P10 2014 OS) \\ \leerfeld $<$ \leerfeld $<$ \leerfeld $<$ \leerfeld
\teil Welche Aussage ist wahr? Kreuze an. \hfill (P10 2015 OS)\\ \kreuz{$2{,}25 < \frac{9}{4}$}\\ \kreuz{$\frac{12}{5} > \frac{9}{4}$}\\ \kreuz{$\sqrt{5} > \frac{9}{4}$}
\teil Welche Aussage ist wahr? Kreuze an. \hfill (P10 2015 OS)\\ \kreuz{$3{,}25 < \frac{13}{4}$}\\ \kreuz{$\frac{17}{5} > \frac{13}{4}$}\\ \kreuz{$\sqrt{10} > \frac{13}{4}$}
% AUSGELASSEN brueche-dezimalzahlen-e5-k2-s9-v7: Missing $ inserted.
% \teil Setze $<$, $=$ oder $>$ ein: $4\,\%$ __ $0{,}4$ \hfill (P10 2018 OS) \\ \leerfeld
\teil \textit{(ausgelassen)}
% AUSGELASSEN brueche-dezimalzahlen-e5-k2-s9-v8: LaTeX Error: Command \item invalid in math m
% \teil Setze $<$, $=$ oder $>$ ein: $0{,}8$ __ $8\,\%$ \hfill (P10 2018 OS) \\ \leerfeld
\teil \textit{(ausgelassen)}
\teil Welche Zahl liegt genau in der Mitte zwischen $-0{,}8$ und $-0{,}7$? \hfill (P10 2020 OS) \\ \leerfeld
\teil Welche Zahl liegt genau in der Mitte zwischen $-1{,}3$ und $-1{,}2$? \hfill (P10 2020 OS) \\ \leerfeld
\end{teile}
\end{aufgabe}
